% ============================================================
% main.tex - 数学建模竞赛论文主文件
% 编译方式: latexmk -xelatex main.tex
% 或在 TeXStudio 中设置编译器为 XeLaTeX
% ============================================================
\documentclass[12pt,a4paper]{ctexart}

% ---- 页面设置 ----
\usepackage[top=2.5cm, bottom=2.5cm, left=2.5cm, right=2.5cm]{geometry}

% ---- 数学 ----
\usepackage{amsmath, amssymb, amsthm}
\usepackage{mathtools}

% ---- 图表 ----
\usepackage{graphicx}
\usepackage{float}
\usepackage{subfigure}
\usepackage{booktabs}       % 三线表
\usepackage{longtable}      % 跨页表格
\usepackage{multirow}
\usepackage{array}

% ---- 代码 ----
\usepackage{listings}
\usepackage{xcolor}
\lstset{
  language=Python,
  basicstyle=\small\ttfamily,
  keywordstyle=\color{blue},
  commentstyle=\color{gray},
  stringstyle=\color{red!70!black},
  numbers=left,
  numberstyle=\tiny\color{gray},
  numbersep=8pt,
  frame=single,
  breaklines=true,
  breakatwhitespace=false,
  tabsize=4,
  showstringspaces=false,
  captionpos=b,
  xleftmargin=2em,
  xrightmargin=1em,
}

% ---- 引用与链接 ----
\usepackage[numbers,sort&compress]{natbib}
\usepackage{hyperref}
\hypersetup{
  colorlinks=true,
  linkcolor=black,
  citecolor=blue!70!black,
  urlcolor=blue!70!black,
  pdfstartview=FitH,
}

% ---- 页眉页脚 ----
\usepackage{fancyhdr}
\pagestyle{fancy}
\fancyhf{}
\fancyhead[C]{\small 数学建模竞赛论文}
\fancyfoot[C]{\thepage}
\renewcommand{\headrulewidth}{0.4pt}

% ---- 标题编号格式 ----
\ctexset{
  section/format = \Large\bfseries\heiti,
  subsection/format = \large\bfseries\heiti,
  subsubsection/format = \normalsize\bfseries\heiti,
}

% ---- 公式编号按章节 ----
\numberwithin{equation}{section}
\numberwithin{figure}{section}
\numberwithin{table}{section}

% ---- 图表路径 ----
\graphicspath{{figures/}}

% ---- 自定义命令 ----
\newcommand{\abs}[1]{\left| #1 \right|}
\newcommand{\norm}[1]{\left\| #1 \right\|}
\newcommand{\RR}{\mathbb{R}}
\newcommand{\NN}{\mathbb{N}}

% ============================================================
\begin{document}

% ---- 标题页 ----
\begin{center}
  {\LARGE\bfseries\heiti 论文标题}\\[1em]
  % 根据竞赛要求可在此添加队号等信息
\end{center}

% ---- 摘要 ----
% ============================================================
% abstract.tex - 摘要与关键词
%
% 硬性要求：摘要页必须独占一整页
% - 页面内只包含"摘要"标题 + 摘要正文 + 关键词
% - 不允许混入目录、正文、页眉等其他内容
% - 摘要正文字数参考比赛规则（国赛 300-500 字，美赛 1 页）
%
% 技术实现：
% - 使用 \thispagestyle{empty} 隐藏摘要页的页眉页脚
% - 用 \vfill 把关键词推到页面底部，视觉上撑满一页
% - 溢出时用 \begin{spacing}{1.25} 或 \small 压缩
% ============================================================

\thispagestyle{empty}

\begin{center}
  {\Large\bfseries\heiti 摘\quad 要}
\end{center}
\vspace{1em}

% ---- 摘要正文 ----
% 结构建议（每段各自独立）：
%   第1段：问题背景 + 整体思路（1-2 句带过）
%   第2段：问题一 —— 方法 + 具体数值结果
%   第3段：问题二 —— 方法 + 具体数值结果
%   第N段：其余问题 ...
%   末段：模型优点 + 局限性 + 推广价值（简短）
%
% 关键：每问必须给出**具体数值结果**，禁止只写"效果良好/结果合理"
%       评委先看摘要，数字就是得分点

本文针对 XXX 问题，建立了 XXX 模型……（示例占位，请替换为实际内容）

针对问题一，采用 XXX 方法，求解得到 XXX = 1256.3（具体数值）……

针对问题二，建立 XXX 模型，结果表明 XXX 比基准方法提升 12.3\%……

（按实际问题数继续段落）

本文模型的优点在于 XXX；不足是 XXX；可推广至 XXX 场景。

\vfill

\noindent\textbf{关键词：} 关键词1\quad 关键词2\quad 关键词3\quad 关键词4


% ---- 目录（按竞赛要求决定是否保留） ----
% 如要求去除目录，注释掉下面两行
\newpage
\tableofcontents
\newpage

% ---- 正文 ----
% ============================================================
% problem_restatement.tex - 问题重述
% ============================================================

\section{问题重述}

% 用自己的语言重述题目背景和问题，切忌照抄原题。
% 可适当增加要研究的问题或添加合理的限制条件。
% 通常 200-300 字。

（问题重述内容）

% ============================================================
% problem_analysis.tex - 问题分析
%
% 本章职责：建立对题目整体的认识，揭示各问之间的逻辑/数据/模型
% 依赖关系。每一问的**详细方法对比**放在 questionX.tex 的第一
% 小节，不要在本章展开。
% ============================================================

\section{问题分析}

\subsection{整体分析}

（用 3-5 句话说明题目的性质、核心难点、解决问题的整体思路。）

\subsection{问题间的依赖关系}

% 按建模阶段产出的依赖图，说明每一问之间是
% 依赖（Q(n+1) 用到 Q(n) 的模型或结果）/
% 扩展（Q(n+1) 在 Q(n) 基础上增加约束或变化条件）/
% 独立（完全独立求解）。
% 可用 tikz 绘制依赖关系图，或用如下的文字 + 表格方式说明。

\begin{table}[htbp]
  \centering
  \caption{问题间的依赖关系}
  \label{tab:dependency}
  \begin{tabular}{clll}
    \toprule
    \textbf{当前问题} & \textbf{依赖的前置问题} & \textbf{关系类型} & \textbf{传递内容} \\
    \midrule
    问题一 & —       & 基线       & —                         \\
    问题二 & 问题一  & 依赖       & 最优决策变量 $x^*$         \\
    问题三 & 问题二  & 扩展       & 增加鲁棒性约束             \\
    问题四 & —       & 独立       & —                         \\
    \bottomrule
  \end{tabular}
\end{table}

\subsection{各问分析概述}

对每一问给出问题类型判断（优化 / 预测 / 评价 / 分类 / 仿真）+ 关键
难点 + 思路方向（1 段即可，详细方法对比放在对应 questionX.tex）。

\textbf{问题一}：（问题类型 + 关键难点 + 思路方向）

\textbf{问题二}：（问题类型 + 关键难点 + 思路方向，如何依赖/扩展问题一）

% 按实际题目数量继续添加

% ============================================================
% assumptions.tex - 模型假设
% ============================================================

\section{模型假设}

% 假设要合理、简化、明确。关键性假设不能缺少。逐条编号。

\begin{enumerate}
  \item 假设数据采集过程中不存在系统误差。
  \item 假设各因素之间相互独立。
  \item （根据实际题目添加）
\end{enumerate}

% ============================================================
% symbols.tex - 符号说明
% ============================================================

\section{符号说明}

% 按论文中出现的顺序罗列。避免一个符号表示多个含义。

\begin{longtable}{clll}
  \toprule
  \textbf{符号} & \textbf{含义} & \textbf{单位} & \textbf{英文名称} \\
  \midrule
  \endfirsthead
  \toprule
  \textbf{符号} & \textbf{含义} & \textbf{单位} & \textbf{英文名称} \\
  \midrule
  \endhead
  \bottomrule
  \endfoot

  $x_i$  & 第$i$个决策变量 & --  & decision variable \\
  $c_j$  & 第$j$个系数     & 元  & coefficient \\
  $Z$    & 目标函数值       & --  & objective value \\
  % 按实际需要继续添加

\end{longtable}


% ---- 各问独立章节 ----
% ============================================================
% question1.tex - 问题一的模型建立与求解
% 每一问使用相同结构，复制此文件为 question2.tex 等
% ============================================================

\section{问题一的模型建立与求解}

% ==============================
\subsection{方法对比与选择}
% ==============================

% 列出至少 2 种可行方法，用对比表展示优劣，
% 正文中对每种方法简要描述原理（各 100-200 字），
% 最后明确说明选择理由。

针对问题一，本文考虑了以下候选方法：

\begin{table}[H]
  \centering
  \caption{问题一候选方法对比}
  \begin{tabular}{lcccc}
    \toprule
    \textbf{方法} & \textbf{适用性} & \textbf{精度} & \textbf{效率} & \textbf{可解释性} \\
    \midrule
    方法A & 高 & 高 & 中 & 高 \\
    方法B & 中 & 高 & 高 & 低 \\
    \bottomrule
  \end{tabular}
\end{table}

（对每种方法的简要描述...）

综合考虑适用性和可解释性，本文选择方法A求解问题一。

% ==============================
\subsection{模型建立}
% ==============================

% 目标函数、约束条件、完整数学表达。

\subsubsection{决策变量}

（变量定义）

\subsubsection{目标函数}

\begin{equation}
  \min Z = \sum_{i=1}^{n} c_i x_i
  \label{eq:q1_obj}
\end{equation}

\subsubsection{约束条件}

\begin{equation}
  \text{s.t.} \quad
  \begin{cases}
    \sum_{i=1}^{n} a_{ij} x_i \leq b_j, & j = 1, 2, \ldots, m \\
    x_i \geq 0, & i = 1, 2, \ldots, n
  \end{cases}
  \label{eq:q1_cons}
\end{equation}

% ==============================
\subsection{求解过程}
% ==============================

% 算法步骤、参数设置、收敛条件。

（求解算法描述、参数设置说明）

% ==============================
\subsection{结果与分析}
% ==============================

% 表格：原始计算数据 / 结果对比
\begin{table}[H]
  \centering
  \caption{问题一求解结果}
  \begin{tabular}{ccc}
    \toprule
    \textbf{指标} & \textbf{数值} & \textbf{单位} \\
    \midrule
    目标函数值 & 1256.3 & 万元 \\
    求解时间   & 2.4    & 秒 \\
    \bottomrule
  \end{tabular}
\end{table}

% 图：至少 4 张，每张至少 100 字分析

\begin{figure}[H]
  \centering
  \includegraphics[width=0.8\textwidth]{q1_fig1.pdf}
  \caption{问题一结果分析图1}
  \label{fig:q1_1}
\end{figure}

如图~\ref{fig:q1_1}所示，（至少100字的详细分析...）

\begin{figure}[H]
  \centering
  \includegraphics[width=0.8\textwidth]{q1_fig2.pdf}
  \caption{问题一结果分析图2}
  \label{fig:q1_2}
\end{figure}

如图~\ref{fig:q1_2}所示，（至少100字的详细分析...）

% 按需添加更多图（每问至少 4 张）

% ---- 灵敏度/稳定性分析 ----
\subsubsection{灵敏度分析}

（对关键参数的灵敏度分析结果）

% ============================================================
% question2.tex - 问题二的模型建立与求解
%
% 完整结构骨架与写法示范见 question1.tex，本文件仅提供占位。
% 后续 question3.tex / question4.tex 复制本文件并替换内容。
%
% 必须遵循的四小节顺序：
%   1. 方法对比与选择  (>= 2 种方法，含对比表 + 各自原理 + 选定理由)
%   2. 模型建立
%   3. 求解过程
%   4. 结果与分析      (>= 4 张图，每张 >= 100 字分析)
% 章节末尾必须加 \FloatBarrier。
% ============================================================

\section{问题二的模型建立与求解}

\subsection{方法对比与选择}

% 候选方法对比表（>= 2 种）+ 每种方法 100-200 字原理 + 选定理由
（内容）

\subsection{模型建立}

% 决策变量 + 目标函数 + 约束条件
（内容）

\subsection{求解过程}

% 算法步骤 + 参数设置 + 收敛条件
（内容）

\subsection{结果与分析}

% 结果表 + >= 4 张图（每图 >= 100 字分析）+ 灵敏度分析
（内容）

\FloatBarrier

% \input{question3.tex}    % 按实际题目数量取消注释
% \input{question4.tex}

% ---- 评价与推广 ----
% ============================================================
% evaluation.tex - 模型评价与推广
% ============================================================

\section{模型评价与推广}

\subsection{模型优点}

% 用整段叙述，不要过度分点。结合具体数据。

（优点描述...）

\subsection{模型不足}

% 诚实指出局限性，结合具体场景和数据。

（不足描述...）

\subsection{模型推广}

% 讨论模型在其他场景的适用性，以及可能的改进方向。
% 不要使用空泛的"前景光明"式结论。

（推广讨论...）


% ---- 参考文献 ----
\newpage
\bibliographystyle{plain}
\bibliography{references}

% ---- 附录 ----
\newpage
% ============================================================
% appendix.tex - 附录
% ============================================================

\appendix
\section*{附录}
\addcontentsline{toc}{section}{附录}

% ==============================
\subsection*{A. 源代码}
% ==============================

% 代码完整原版粘贴，使用 lstlisting 格式。
% 按问题编号分节。

\subsubsection*{A.1 问题一求解代码}

\begin{lstlisting}[language=Python, caption={问题一求解代码 (question1\_solve.py)}]
# -*- coding: utf-8 -*-
"""
问题一求解代码
"""
import numpy as np
# ... 完整代码粘贴于此 ...
\end{lstlisting}

\subsubsection*{A.2 问题二求解代码}

\begin{lstlisting}[language=Python, caption={问题二求解代码 (question2\_solve.py)}]
# ... 完整代码粘贴于此 ...
\end{lstlisting}

% ==============================
\subsection*{B. 运行数据与补充表格}
% ==============================

% 正文中未展示的详细数据表格、中间计算结果等。
% 使用 longtable 处理大表格。

\begin{longtable}{cccc}
  \caption{补充数据表} \\
  \toprule
  \textbf{序号} & \textbf{数据1} & \textbf{数据2} & \textbf{数据3} \\
  \midrule
  \endfirsthead
  \caption{补充数据表（续）} \\
  \toprule
  \textbf{序号} & \textbf{数据1} & \textbf{数据2} & \textbf{数据3} \\
  \midrule
  \endhead
  \bottomrule
  \endfoot

  1 & -- & -- & -- \\
  % ...
\end{longtable}


\end{document}
